\scanpage{56}
fibrés de $=$ espèce sur $S$, sous réserve que
$T\to S$ soit homogène.

En particulier, si $G$ est un groupe fini
ordinaire, un \deleted{\unsure{revêtem.}} revêtement principal de
groupe $G$ de $Y/T$ est transformé en
un \deleted{\unsure{revêtement}} fibré principal \unsure{hom|hom.|homogène} sur
$f_\ast(Y/T)$, de groupe $f_\ast(Y/T)\times_S f_\ast(\deleted{\unsure{G\times S}}T\times G/T$.
% slanted
\margin{\unsure{Pas|Pour} la peine de supposer $T/S$ fini, \unsure{pourvu que l'on sache|faisons que l'on sache} que la \unsure{descente} est possible.}

Cela montre l'intérêt de déterminer, pour
un ens fini $I$ \deleted{\unsure{variable|ordinaire}}, la variété
\deleted{\unsure{un}} étale $f_\ast(T\times I/T)=\varphi(I)$\deleted{\unsure{C'est|qui est telle que}}

\deleted{\unsure{$\varphi$} On voit que $\varphi(T)(I)$ est un foncteur
\deleted{\unsure{fini}} \uline{multiplicatif} \inserted{et $=$ exact à gauche} qui va des ens.
finis aux variétés étales de $S$. Choisissons
un pt géométrique $a$ sur $S$ supposé
connexe, on obtient ainsi un
foncteur $I\mapsto\varphi_a(T)\deleted{(I)}$ (Ens finis) dans les
ens \inserted{finis} où $\pi_1(S,a)$ opère. En tant que
foncteur à valeurs dans les ens finis,
il est de la forme $\operatorname{Hom}_{\mathrm{Ens}\,f}(A_a(T),I)$,
[petit raisonnement], et comme
$\pi_1(S,a)$ opère dans ce foncteur, il opère
dans $A_a(T)$, \deleted{qui} on a ainsi défini}
% slanted
\margin{\unsure{prendre pour principe un|pour plus de simplicité}}
% archive sheet number: 55
